\chapter{Introducing \lupnt{}}
\label{chap:introduction}

This chapter states what \lupnt{} is, how the C++ library and the \pylupnt{}
Python package fit together, and where every model documented in the rest of
this manual lives in the source tree. It is the map: each module named here is
specified in full in a later chapter. The library was first described in
\citep{iiyama2023lupnt}, extended to communications and the agent-based
simulation framework in \citep{casadesus2025lupnt}, and is used as the
computational backbone of \citep{iiyama2026thesis}.

\section{What \lupnt{} is}
\label{sec:intro-what}

\lupnt{} is an open-source C++/Python library for Lunar Positioning,
Navigation, and Timing (PNT) research. It provides high-fidelity
astrodynamics, signal-propagation models, measurement models, and navigation
algorithms tailored for cislunar missions. It is released under the MIT
License.

Although the library is designed for lunar PNT analysis, nothing in its core is
lunar-specific. The central body of a dynamics model, a frame conversion, or a
constellation is a parameter, so the same machinery covers spacecraft orbit
determination and navigation research for other mission scenarios --- from
Earth orbit (LEO and GNSS) to interplanetary (Mars) --- as several of the
bundled examples demonstrate.

Three properties distinguish \lupnt{} from a general-purpose astrodynamics
toolkit:

\begin{description}[leftmargin=1.2cm,style=nextline,font=\normalfont\itshape]
  \item[Differentiable by construction] Every state, dynamics, and measurement
    function is templated on the scalar type \cpp{Real}, which is an
    \code{autodiff} forward-mode dual number \citep{leal2018autodiff,
    griewank2008}. Jacobians used by the filters are therefore machine-exact
    rather than finite-differenced.
  \item[Relativistically explicit] Time is not a floating-point number of
    seconds since an arbitrary origin. \lupnt{} carries an explicit time scale
    with every epoch --- $\TAI$, $\TT$, $\TCG$, $\TCB$, $\TDB$, $\TCL$, $\LT$,
    $\UTC$, $\UTone$, $\GPST$ --- including the lunar scales needed for
    cislunar timekeeping \citep{turyshev2025transformations,
    turyshev2026cislunar}.
  \item[Config-driven] Scenarios are described in YAML rather than in code, so
    a new mission study is a configuration file plus, at most, one new
    \cpp{Application} subclass.
\end{description}

\begin{lupntnote}
\lupnt{} is installed \emph{from source only}; it is no longer distributed on
PyPI. The legacy \pylupnt{} wheels (last published for version \code{0.0.4})
predate the current C++/Fortran dependency stack --- OpenCV, Matplot++, the
Fortran plasma models, HDF5 --- and are unmaintained. \cref{chap:getting-started}
is the supported installation path.
\end{lupntnote}

\section{Architecture: agents, applications, simulations}
\label{sec:intro-architecture}

\lupnt{} is organised as a \emph{config-driven, agent-based simulation
framework}. Three abstractions carry the whole design:

\begin{description}[leftmargin=1.2cm,style=nextline,font=\normalfont\itshape]
  \item[Agent] A physical entity with a state and a propagator: a satellite, a
    spacecraft, a ground station, a rover, a lander, a surface station, or a
    whole constellation. An agent owns its \emph{dynamics} (how its state
    evolves) and its \emph{devices} (clock, IMU, camera, GNSS receiver,
    communication terminals). Implemented in \file{cpp/lupnt/agents}.
  \item[Application] The mission and navigation logic hosted \emph{by} an
    agent: an orbit-determination and time-synchronisation (ODTS) filter, an
    inter-satellite-link (ISL) estimator, a surface-rover or lander navigation
    filter, an ephemeris or almanac generator. Applications subscribe to
    scheduled events and produce measurements, state updates, and logs.
    Implemented in \file{cpp/lupnt/applications}.
  \item[Simulation] The event-scheduled engine that owns the world clock, the
    agent registry, and the event queue, and that steps every agent and
    application forward in a consistent time scale. Implemented in
    \file{cpp/lupnt/simulations/simulation.cc} and
    \file{cpp/lupnt/simulations/world.cc}.
\end{description}

Agents, their devices, their dynamics, and their applications are all
assembled from YAML by an \emph{asset factory}
(\file{cpp/lupnt/core/asset_factory.h}), so a new scenario is described in
configuration rather than code. The reusable configuration fragments ship in
\file{configs}; the end-to-end scenarios currently provided there are
\file{configs/ephemeris.yaml}, \file{configs/ground_station_odts.yaml},
\file{configs/isl_odts_distributed.yaml}, \file{configs/lander_nav.yaml},
\file{configs/lunar_gnss_odts.yaml}, \file{configs/sat_bearing_odts.yaml}, and
\file{configs/surface_rover_nav.yaml}, all layered on the shared
\file{configs/config.yaml}. \cref{chap:new-simulation} walks through the
architecture and the workflow for building a scenario end to end.

\section{Module map}
\label{sec:intro-modules}

\cref{tab:intro-modules} is the functional view of the library: what each
module does. \cref{tab:intro-dirs} is the physical view: which directory under
\file{cpp/lupnt} implements it and which chapter of this manual specifies it.

\begin{table}[H]
  \centering
  \caption{Functional module map of \lupnt{}.}
  \label{tab:intro-modules}
  \begin{tabularx}{\textwidth}{@{}lL@{}}
    \toprule
    Module & Description \\
    \midrule
    Agents \& Applications &
      Config-driven agents (satellite, ground station, rover, lander, surface
      station, constellation) that host navigation and mission
      \emph{Applications} --- ODTS, ISL, surface-rover and lander navigation,
      ephemeris and almanac generation --- built from YAML via an asset
      factory. \\
    Simulations &
      Event-scheduled \cpp{Simulation} engine plus end-to-end scenario drivers
      (ground-station and inter-satellite ODTS, lunar GNSS ODTS, lander and
      surface navigation, ephemeris fitting). \\
    Devices \& sensors &
      Clocks, IMUs, cameras, GNSS receivers, and communication devices attached
      to agents. \\
    Dynamics &
      Two-body, $N$-body, and high-fidelity force models (gravity, drag, solar
      radiation pressure) for Earth, lunar, and arbitrary central-body orbits. \\
    Environment &
      Gravity fields, atmosphere, solar-system bodies, occultation, and
      plasma/ionosphere models. \\
    Plasma &
      GCPM v2.4 plasmasphere \citep{gallagher2000gcpm} plus IRI ionosphere
      \citep{bilitza2022iri} electron density; GNSS ray tracing with TEC and
      signal delay for cislunar links (the integrated \code{pecsim} code). \\
    GNSS &
      GPS/GNSS signal generation, SP3- and ANTEX-backed constellations,
      yaw-steering and attitude models, and space-user ranging. \\
    Measurements &
      Pseudorange, Doppler, and carrier-phase measurement models with
      light-time, Shapiro, and optional plasma corrections. \\
    Filters \& States &
      Extended Kalman filter (EKF), unscented Kalman filter (UKF), square-root
      information filter (SRIF) and smoother, $\mat{U}\mat{D}\mat{U}\transpose$
      and Schmidt
      variants, and batch least-squares estimators over composable state and
      joint-state abstractions. \\
    Interfaces &
      Data loaders and I/O: SP3, ANTEX, RINEX navigation, TLE, SPICE kernels,
      EOP and TAI--UTC tables, LOLA DEM and crater data, plus Cesium and
      Matplotlib export. \\
    Conversions &
      Reference-frame transformations (ECI, ECEF, LVLH, Moon-centred, generic
      body-fixed and inertial) and time-system utilities. \\
    Numerics &
      Numerical integration (RK4, RK8, RKF45), a Nelder--Mead optimiser,
      Chebyshev fitting, interpolation, and matrix utilities. \\
    Visualization &
      Matplotlib and Plotly plotting plus interactive 3-D CesiumJS
      (\url{https://cesium.com/platform/cesiumjs/}) scenes of constellations
      and surface assets (\py{pnt.plot.CesiumScene}). \\
    Python bindings &
      Full \pylupnt{} package exposing the C++ library via pybind11 ---
      including subclassing \cpp{Application}, \cpp{Agent}, and
      \cpp{Measurement} in pure Python and registering them into a YAML-driven
      \cpp{Simulation}. \\
    \bottomrule
  \end{tabularx}
\end{table}

\begin{table}[H]
  \centering
  \caption{Source directories under \code{cpp/lupnt} and where they are
           specified in this manual.}
  \label{tab:intro-dirs}
  \begin{tabularx}{\textwidth}{@{}llL@{}}
    \toprule
    Directory & Specified in & Principal contents \\
    \midrule
    \file{agents/} & \cref{chap:new-simulation} &
      \cpp{Agent} base class; \cpp{Satellite}, \cpp{Spacecraft},
      \cpp{GroundStation}, \cpp{Rover}, \cpp{Lander}, \cpp{SurfaceStation};
      \cpp{Constellation}, \cpp{GnssConstellation},
      \cpp{LunarNavConstellation}; the \cpp{EphemerisManager},
      \cpp{GroundStationManager}, and \cpp{SurfaceStationManager}
      registries. \\
    \file{applications/} & \cref{chap:new-simulation} &
      \cpp{Application} base plus per-mission subdirectories
      \file{ephemeris/}, \file{ground_station/}, \file{lander/},
      \file{lunar_gnss_odts/}, \file{lunar_sat_odts/},
      \file{lunar_station/}, \file{rover/}. \\
    \file{simulations/} & \cref{chap:new-simulation} &
      \file{simulation.cc}, \file{world.cc}, and the \file{ephemeris/} and
      \file{lunar_gnss_odts/} scenario drivers. \\
    \file{devices/} & \cref{chap:measurements} &
      \file{device.h} base, \file{clock.cc}, \file{imu.cc}, \file{camera.cc},
      \file{gnss_device.cc}, \file{comm_devices.cc}. \\
    \file{dynamics/} & \cref{chap:dynamics} &
      \file{numerical_orbit_dynamics.cc}, \file{analytical_orbit_dynamics.cc},
      \file{attitude_dynamics.cc}, \file{clock_dynamics.cc},
      \file{imu_dynamics.cc}, \file{joint_orbit_clock_dynamics.cc},
      \file{parameter_dynamics.h}, \file{surface_dynamics.cc},
      \file{dynamics_params.cc}. \\
    \file{environment/} & \cref{chap:dynamics} &
      \file{body.cc}, \file{forces.cc}, \file{gravity.h},
      \file{atmosphere.cc}, \file{occultation.cc}, \file{solar_system.cc};
      the \file{plasma/} subtree (\file{core/}, \file{env/}, \file{fortran/},
      \file{gcpm/}, \file{igrf/}, \file{tec/}). \\
    \file{measurements/} & \cref{chap:gnss-measurements},
      \cref{chap:measurements} &
      \file{measurement.h}, \file{gnss_measurement.cc},
      \file{crosslink_measurement.cc}, \file{ground_range_measurement.cc},
      \file{ground_station_corrections.cc}, \file{lander_measurements.cc},
      \file{lunar_gnss_combined_measurement.cc},
      \file{surface_measurements.cc}, \file{antenna.cc}, \file{channel.cc},
      \file{comms_utils.cc}. \\
    \file{attitude/} & \cref{chap:gnss-measurements} &
      \file{gnss_attitude.cc}, \file{gnss_yaw_steering.cc}
      \citep{kouba2009yaw, cheng2025yaw}. \\
    \file{transmission/} & \cref{chap:link-budget} &
      \file{transmission.cc} --- link geometry and signal transmission
      modelling. \\
    \file{states/} & \cref{chap:filters} &
      \file{state.cc}, \file{joint_state.cc}, \file{params.h}. \\
    \file{interfaces/} & \cref{chap:sp3-download} &
      \file{sp3_loader.cc}, \file{antex_loader.cc},
      \file{rinex_nav_loader.cc}, \file{tle.cc}, \file{eop.cc},
      \file{tai_utc.cc}, \file{spice.cc}, \file{spice_cheby.cc},
      \file{kernels.cc}, \file{iau_sofa.cc}, \file{dem.cc},
      \file{lola_dem.cc}, \file{crater_data.h}, \file{cesium.cc},
      \file{matplot.h}, \file{python.cc}, \file{yaml.h}. \\
    \file{conversions/} & \cref{chap:time-conversions},
      \cref{chap:frame-conversions} &
      \file{time_conversions.cc}, \file{epoch.cc},
      \file{frame_conversions.cc}, \file{frame_converter.cc},
      \file{frame_converter_spice.cc}, \file{coordinate_conversions.cc},
      \file{state_conversions.cc}, \file{state_converter.cc},
      \file{anomaly_conversions.cc}, \file{attitude_conversions.cc},
      \file{mean_osc_conversions.cc}. \\
    \file{numerics/} & \cref{chap:integration}, \cref{chap:filters} &
      \file{integrator.cc}, \file{interpolation.cc}, \file{math_utils.cc},
      \file{cheby_fit.h}, \file{graphs.cc}, \file{vector_macros.h}; the
      \file{filters/} subtree (\file{ekf.cc}, \file{ukf.cc}, \file{srif.cc},
      \file{udu_filter.cc}, \file{schmidt_ekf.h}, \file{batch_filter.cc},
      \file{adaptive_process_noise.cc}). \\
    \file{core/} & \cref{chap:new-simulation} &
      \file{constants.h}, \file{config.cc}, \file{definitions.cc},
      \file{asset_factory.h}, \file{data_logger.cc}, \file{event.h},
      \file{file.cc}, \file{logger.cc}, \file{object.h},
      \file{progress_bar.cc}, \file{random_engine.cc},
      \file{string_utils.cc}, \file{error.h}. \\
    \bottomrule
  \end{tabularx}
\end{table}

The single umbrella header \file{cpp/lupnt/lupnt.h} is generated at configure
time by \file{cmake/WriteHeader.cmake} and includes every public header, so a
consumer only ever writes \code{\#include <lupnt/lupnt.h>}.

\section{The C++ core and the Python layer}
\label{sec:intro-bindings}

\lupnt{} is a C++20 library with a Fortran component (the legacy IRI, GCPM,
IGRF, and \code{xform} plasma codes, compiled with
\code{-ffixed-form -std=legacy -fno-automatic}). The CMake project therefore
declares \code{LANGUAGES CXX Fortran}. Everything else is a thin layer on top:

\begin{enumerate}[leftmargin=*]
  \item \file{cpp/lupnt} builds the shared library target \code{lupnt}. Sources
    are globbed (\code{file(GLOB\_RECURSE ...\ CONFIGURE\_DEPENDS)}), so adding
    a \file{.cc} file requires no CMake edit.
  \item \file{python/bindings} is a \emph{separate} CMake project that links
    against \code{lupnt} and produces the pybind11 extension module
    \code{\_pylupnt}. Its \code{pylupnt-dev} target copies the built shared
    object into \file{python/pylupnt/} so the package is usable in place, with
    no install step.
  \item \file{python/pylupnt} is the Python package. \file{__init__.py}
    re-exports the entire extension module (\code{from .\_pylupnt import *}),
    so \py{pylupnt} symbols and C++ symbols share one namespace. The
    pure-Python sub-packages \file{python/pylupnt/plot},
    \file{python/pylupnt/plasma}, \file{python/pylupnt/core}, and
    \file{python/pylupnt/interfaces} add plotting, Cesium export, plasma
    helpers, and data-path management. A generated stub file
    \file{python/pylupnt/_pylupnt.pyi} gives editors and \code{mypy} full
    type information for the bound API.
\end{enumerate}

The binding layer is not read-only: \cpp{Application}, \cpp{Agent}, and
\cpp{Measurement} are exposed with pybind11 trampolines, so they can be
subclassed in pure Python and registered into a YAML-driven \cpp{Simulation}.
The canonical demonstration is example~17
(\file{python/examples/ex17_python_new_sim_example.ipynb}), an angles-only
orbit-determination scenario whose application and measurement model are
written entirely in Python.

\subsection{External dependencies}
\label{sec:intro-deps}

Dependencies come from two places. Conda-forge packages are resolved by Pixi
(\cref{sec:gs-install}) and found with \code{find\_package}: Eigen~3.4
\citep{eigen2010}, \code{fmt}, \code{pybind11}, \code{yaml-cpp},
\code{spdlog}, OpenCV, GDAL, OpenMP, HDF5, TBB, Boost, and \code{gfortran}.
Header- and source-level dependencies are vendored at configure time by CPM
(\file{cpp/thirdparty/CMakeLists.txt}), as listed in
\cref{tab:intro-thirdparty}.

\begin{table}[H]
  \centering
  \caption{Dependencies fetched by CPM at configure time.}
  \label{tab:intro-thirdparty}
  \begin{tabularx}{\textwidth}{@{}llL@{}}
    \toprule
    Package & Version & Role \\
    \midrule
    \code{autodiff}          & \code{v1.1.2}          & forward-mode dual numbers behind the \cpp{Real} scalar type \citep{leal2018autodiff} \\
    \code{cspice}            & \code{master}          & NAIF SPICE toolkit: kernels, ephemerides, body-fixed frames \citep{acton1996spice} \\
    \code{matplotplusplus}   & \code{v1.2.2}          & C++ plotting used by the C++ examples \\
    \code{magic\_enum}       & \code{v0.9.7}          & compile-time enum reflection, used by the YAML asset factory \\
    \code{HighFive}          & \code{v3.0.0-beta2}    & HDF5 C++ interface for the data logger \\
    \code{crow}              & \code{v1.2.1}          & embedded HTTP server for the Cesium scene viewer \\
    \code{PackageProject.cmake} & \code{v1.11.2}      & generates the installable \code{lupntConfig.cmake} package \\
    \bottomrule
  \end{tabularx}
\end{table}

\section{Repository layout}
\label{sec:intro-layout}

\begin{lstlisting}[language=shell, caption={Top-level layout of the repository.}, label={lst:intro-layout}]
lupnt/
  cpp/
    lupnt/            # C++ library source, one directory per module
    examples/         # C++ example programs
      tutorials/      # C++ counterparts to the Python notebooks (ex1-ex16)
    test/             # Catch2 / CTest suites, plus orekit/ and gmat/ fixtures
    thirdparty/       # CPM declarations for autodiff, cspice, matplot++, ...
  python/
    pylupnt/          # Python package, installed in place
      core/           # data-path and download helpers
      plasma/         # pylupnt.plasma sub-module
      plot/           # Matplotlib / Plotly / Cesium plotting
      interfaces/     # Python-side data interfaces
      _pylupnt*.so    # compiled pybind11 extension (built by `pixi run build-py`)
      _pylupnt.pyi    # type stubs for the bound API
    bindings/         # pybind11 binding sources (separate CMake project)
    examples/         # tutorial notebooks ex1-ex17
    test/             # pytest suite
  configs/            # YAML scenario configuration
  projects/           # research workflows, notebooks, and scripts
  data/
    LuPNT_data/       # runtime data, fetched automatically on the first build
      antenna/  ephemeris/  examples/  gnss/  gravity/
      ground_station/  planetary_coeff/  plasma/  surface/  tle/  topo/
  scripts/            # install_kernel.py, notebook runners, format checks
  cmake/              # CPM.cmake, FetchLuPNTData.cmake, WriteHeader.cmake, ...
  docs/               # Sphinx / Breathe / Exhale sources (API reference site)
  latex_docs/         # this manual: lupnt.tex, preamble.tex, chapters/, refs.bib
  pixi.toml           # reproducible environment and task definitions
\end{lstlisting}

The bundled runtime dataset \file{data/LuPNT_data} is gitignored and roughly
650~MB; it holds SPICE kernels and planetary ephemerides
\citep{park2021de440}, spherical-harmonic gravity fields
\citep{lemoine2013grgm}, GNSS antenna and clock products, EOP and TAI--UTC
tables \citep{iers2010}, TLEs, plasma coefficient files, and lunar surface
imagery and DEMs. It is downloaded automatically the first time the project is
built; see \cref{sec:gs-build}.

\section{Worked examples}
\label{sec:intro-examples}

A progressive set of Jupyter notebooks under \file{python/examples} teaches the
\pylupnt{} API and reproduces \lupnt{}'s core navigation workflows, from a
single orbit propagation up to full orbit-determination filters, constellation
design, optical navigation, and non-lunar PNT. Each has a C++ counterpart under
\file{cpp/examples/tutorials}. \cref{chap:tutorials} indexes them and maps each
one to the models it exercises; the short version is
\cref{tab:intro-examples}.

\begin{table}[H]
  \centering
  \caption{The bundled tutorial examples.}
  \label{tab:intro-examples}
  \begin{tabularx}{\textwidth}{@{}clL@{}}
    \toprule
    \# & Example & Theme \\
    \midrule
     1 & \file{ex1_propagate_orbit}        & elements, frames, force models \\
     2 & \file{ex2_time_conversions}       & $\TT$, $\TCG$, $\TCB$, $\TDB$, $\TCL$, $\LT$ \\
     3 & \file{ex3_gnss_interface}         & TLEs, antenna gain, SP3 versus broadcast \\
     4 & \file{ex4_plasmasphere}           & GCPM v2.4 electron density, TEC, ray-traced delay \\
     5 & \file{ex5_gnss_measurement_sim}   & lunar-receiver visibility and $C/N_0$ \\
     6 & \file{ex6_gnss_odts}              & weak-signal EKF for position, velocity, clock, SRP \\
     7 & \file{ex7_groundstation_odts}     & DSN batch least squares plus SRIF/smoother \\
     8 & \file{ex8_isl_odts}               & five parallel consider-state EKFs, surface timing \\
     9 & \file{ex9_ephemeris}              & compressing trajectories into broadcast models \\
    10 & \file{ex10_surface_rover}         & IMU, LCRNS pseudoranges, DEM error-state EKF \\
    11 & \file{ex11_lander_navigation}     & powered-descent MEKF: IMU, altimeter, craters \\
    12 & \file{ex12_constellation_design}  & ELFO Walker layout, PDOP, coverage, EIRP \\
    13 & \file{ex13_cesium}                & interactive 3-D CesiumJS scenes \\
    14 & \file{ex14_opnav}                 & lunar-horizon image processing to EKF fixes \\
    15 & \file{ex15_marspnt}               & Mars gravity, frames, Walker constellation \\
    16 & \file{ex16_leopnt}                & Earth LEO with Harris--Priester drag \citep{harris1962priester} \\
    17 & \file{ex17_python_new_sim_example} & authoring an Application and Measurement in Python \\
    \bottomrule
  \end{tabularx}
\end{table}

\section{Quick start}
\label{sec:intro-quickstart}

All dependencies --- compiler toolchain, CMake, Ninja, Python, and every
C++/Python package --- are provisioned by Pixi (\url{https://pixi.sh}) from
conda-forge, so no compiler or library is installed by hand. From the
repository root:

\begin{lstlisting}[language=shell]
# repository root
pixi install         # provision the conda-forge environment into .pixi/
pixi run build       # build the C++ library (target: lupnt)
pixi run build-py    # build and deploy the Python bindings (_pylupnt)
\end{lstlisting}

The runtime dataset is downloaded on the first build.
\cref{chap:getting-started} gives the complete procedure, including the NASA
Earthdata credentials, the Windows/WSL path, and how to verify the install; \cref{chap:development} covers build variants, debugging, and the
VS~Code workflow; \cref{chap:building-docs} covers building the documentation.

\section{How to read this manual}
\label{sec:intro-howtoread}

This manual is a specification, not a tour. Read it as follows.

\begin{itemize}[leftmargin=*]
  \item If you want a running install, read \cref{chap:getting-started} and
    stop there. Everything else assumes a working \code{pixi} environment.
  \item If you want to build a scenario, read \cref{chap:new-simulation} for
    the agent/application/simulation contract, then \cref{chap:tutorials} for
    the closest worked example, and copy it.
  \item If you want to know \emph{what a model actually computes} --- its
    state and unit contract, its governing equations with every symbol
    defined, the code that implements them, the configuration knobs that
    change its behaviour, and the explicit statement of what it approximates
    and omits --- go to \cref{part:math}. That part is the core of this
    document. Its chapters are ordered so that each one depends only on the
    ones before it: time (\cref{chap:time-conversions}), frames
    (\cref{chap:frame-conversions}), dynamics (\cref{chap:dynamics}),
    integration and differentiation (\cref{chap:integration},
    \cref{chap:autodiff}), estimation (\cref{chap:filters}), then the
    measurement and signal models (\cref{chap:gnss-measurements} onward).
  \item If you are chasing a number rather than a method, \cref{part:verification}
    holds the external data products, the cross-validation results against
    GMAT \citep{gmat2023} and Orekit, and the publication list.
\end{itemize}

Throughout, \code{monospace} marks an identifier, an enum value, or a command,
and \file{a/slashed/path} is a path relative to the repository root. Every
listing names the source file it came from on its first line.

\section{Citing \lupnt{}}
\label{sec:intro-citing}

Work that uses \lupnt{} should cite the original simulator paper
\citep{iiyama2023lupnt} and, for the current agent-based communications and PNT
framework, \citep{casadesus2025lupnt}. The full development history, the
derivations behind the models in \cref{part:math}, and the mission-design
studies built on them are collected in \citep{iiyama2026thesis}.

\begin{lstlisting}[language=bibtexlang]
@inproceedings{IiyamaCasadesus2023,
  title     = {LuPNT: Open-Source Simulator for Lunar Positioning,
               Navigation, and Timing},
  author    = {Iiyama, Keidai and Casadesus Vila, Guillem and Gao, Grace},
  booktitle = {Proceedings of the Institute of Navigation GNSS+ Conference
               (ION GNSS+ 2023)},
  institution = {Stanford University},
  year      = {2023},
  url       = {https://github.com/Stanford-NavLab/LuPNT},
}
\end{lstlisting}

The MIT License covers \lupnt{}'s own source code only. The plasma module
embeds third-party scientific codes that retain their original terms and
carry their own citation requirements: IRI2007 and IRI2020
\citep{bilitza2022iri}, GCPM v2.4 \citep{gallagher2000gcpm}, the \code{xform}
coordinate-transformation utilities, and IGRF-14 \citep{alken2021igrf}.
